In this paper, we improve the above results from~\cite{CKNV16} for bipartite regular graphs. Let $b(k,\theta)$ denote the maximum order of a connected bipartite $k$-regular graph $\Gamma$ with $\lambda_2(\Gamma) \leq \theta$. Bipartite regular graphs $\Gamma$ with $\lambda_2(\Gamma) \leq \theta$ have been classified for $\theta=\sqrt{2}$~\cite{TY00}, $\theta=\sqrt{3}$~\cite{KS13}, and $\theta=2$~\cite{KS14}. We obtain a general upper bound for $b(k,\theta)$ for any $0\leq \theta< 2\sqrt{k-1}$. Our bound gives the exact value of $b(k,\theta)$ whenever there exists a bipartite distance-regular graph of degree $k$ with second largest eigenvalue $\theta$, diameter $d$ and girth $g$ such that $g\geq 2d-2$. For certain values of $d$, there are infinitely many such graphs of various valencies $k$. When $d=11$ or $d\geq 15$, we prove the non-existence of bipartite distance-regular graphs with $g\geq 2d-2$. Our results generalize previous work of H\o holdt and Justesen~\cite{HJ14} obtained in their study of graph codes and imply some results of Li and Sol\'{e}~\cite{CLS96} relating the second largest eigenvalue of a bipartite regular graph to its girth. The degree-diameter or Moore problem for graphs~\cite{MS} is about determining the largest graphs of given maximum degree and diameter. Given the connections between the diameter and the second largest eigenvalue of bipartite regular graphs (see~\cite{Ch89} for example), our Theorem~\ref{bound} can be interpreted as a spectral version of the Moore problem for bipartite regular graphs. 
